\section{Ergebnisse}
\subsection{Bequemlichkeit und Effizienz}
Ein zentrales Motiv, welches sich aus allen drei Interviews ergab, ist die wahrgenommene Bequemlichkeit des Online-Shoppings im Vergleich zum stationären Einzelhandel. Hier beschrieben die Interviewpartner wiederholt, dass Online-Einkäufe vor allem aufgrund der Schnelligkeit und dem geringeren Aufwand bevorzugt wurden.

Ein Beispiel, welches sich aus einem Interview ergab lautet:\\  \enquote{Sachen oder Dinge, die schnell gehen mussten, vielleicht online.} (Interview 1).\\  Auch im zweiten Interview wurde die niedriegere Hürde betont:\\ \enquote{Dann geht das einfach schneller: Ach, ich brauche dies, ich brauche das, mal kurz … angeklickt, bestellt.} (Interview 2).

Es wurde durch die Befragten mehrfach darauf hingewiesen, dass die Corona-Pandemie bestehende Bequemlichkeiten nur noch weiter verstärkt hat. Hier zeigte sich, dass Convenience ein zentraler Treiber für die Verlagung vom Einzelhandel zum Online-Shopping war.

\subsection{Vergleichbarkeit und erweiterete Informationsmöglichkeiten}
Ein weiteres Motiv neben der Bequemlichkeit war die verbesserte Vergleichbarkeit von Produkten als wesentlicher Vortiel des Online-Shoppings. Die Interviewpartner beschrieben hier insbesonder die Möglichkeit Preise, Bewertungen und Produktmerkmale schnell gegegüberzustellen.

Ein konkretes Beispiel aus dem ersten Interview lautet:\\  \enquote{Vorteil ist natürlich, dass man im Onlineshop besser vergleichen kann, Preise vergleichen kann oder wo etwas verfügbar ist.}.\\  Im zweiten Interview wurde zusätzlich noch die dazu gewonnene Informationssuche betont:\\ \enquote{Ich habe mir das immer angeguckt und dann habe ich einfach mehr Testberichte gelesen.}.\\  Auch der dritte Interviewpartner vierwies auf die Bedeutung von Bewertungen und Rezensionen:\\  \enquote{während ich online auf viele Bewertungen zugreifen kann und Vergleichsseiten nutzen kann.}.

Diese Ergebnisse zeigen, dass das Online-Shopping nicht nur als Alternative, sondern auch als wichtiges Recherchemedium wahrgenommen wird.

\subsection{Gewohnheitsbildung durch die Pandemie}
Ein auch zentrales, wiedekehrendes Muster war die nachhaltige Veränderung des Kaufverhalten. Mehrere Interviewpartner berichteten, dass sie während der Corona-Pandemie neue Gewohnheiten etablierten, die sie teilweise bis heute weiter beibehielten.

Hier berichtete die befragte Person aus Interview 1:\\  \enquote{Ja, beim Online-Shopping sind Gewohnheiten geblieben.}.\\ Dies konnte auch durch die befragte Person aus dem dritten Interview bestätigt werden:
\\  \enquote{Während der Corona-Pandemie habe ich noch mehr online eingekauft, weil viele Geschäfte geschlossen hatten.}.

Es wurde allerdings auch gezeigt, dass die Entwicklung dieser Gewohntheiten nicht vollständig linear ablief und teilweise eine bewusste Rückkehr zum stationären Handel beschrieben wurde.